\documentclass[11pt,a4paper]{letter}
\usepackage[utf8]{inputenc}
\usepackage[T1]{fontenc}
\usepackage[french]{babel}
\usepackage[left=2.5cm,right=2.5cm,top=2cm,bottom=2cm]{geometry}
\usepackage{xcolor}

% --- Couleurs Eviden ---
\definecolor{primary}{RGB}{0, 55, 110}

\signature{Loïc Aron MBASSI EWOLO}
\address{Loïc Aron MBASSI EWOLO\\
Cergy, France\\
(+33) 7 59 52 33 98\\
loicaron.mbassiewolo@ensea.fr\\
LinkedIn : loic-mbassi | GitHub : Nameless0l}

\begin{document}

\begin{letter}{Eviden / Atos\\
Pôle Cellulaire Avantix\\
Aix-en-Provence}

\opening{Madame, Monsieur,}

Élève-ingénieur en 2A à l'ENSEA (double diplôme avec Polytechnique Yaoundé), je candidate au stage de développement GPU / traitement du signal.

Je suis à l'aise avec \textbf{Python}, \textbf{NumPy} et \textbf{PyTorch} -- c'est ce que j'ai utilisé au MILA cet été pour analyser la convergence de modèles sur le papier Grokking d'OpenAI. J'ai aussi une base solide en \textbf{C/C++} (projet système multiprocessus, embedded sur STM32).

Côté GPU, j'ai travaillé sur \textbf{Nvidia Jetson} lors du challenge CoHoMa pour déployer du YOLOv10 en temps réel. Je n'ai pas encore codé directement en CUDA, mais je comprends le modèle d'exécution et je suis motivé pour apprendre.

Pour le \textbf{traitement du signal}, je suis les cours à l'ENSEA (filtrage, FFT, analyse spectrale) et j'ai travaillé sur du traitement de signaux ECG et de données capteurs IMU. Les maths derrière (algèbre linéaire, tenseurs) ne me font pas peur -- c'était mon domaine à Polytechnique.

Je travaille sous \textbf{Linux} au quotidien et je documente rigoureusement mon travail (habitude prise au MILA).

Ce stage m'intéresse pour deux raisons : approfondir le calcul GPU sur du matériel récent (DGX Spark / RTX5090), et voir comment ça s'applique au traitement du signal en conditions réelles. La possibilité d'embauche ensuite est aussi un plus.

Disponible pour un entretien.

\closing{Cordialement,}

\end{letter}
\end{document}
